\section{Arity-normalized staged proof compilation}
\label{sec:compiler}

This section contains the central contribution.  The mathematical graph machinery in the existing packets
is sound in spirit; the persistent engineering obstacle is turning local theorems with different free
variables and contexts into one final \(\Prf\,\PRAAx\) term.  The solution uses only three structural
theorems already proved in \texttt{Deduction.lean}: substitution, weakening, and simultaneous cut.

\subsection{Local theorems and a global parameter frame}

A domain library naturally proves lemmas in small local frames.  For example,
\begin{itemize}
  \item a Boolean comparison lemma may have arity two;
  \item a matrix multiplication lemma may have arity containing dimension, two matrices, and an index;
  \item a trace step lemma may bind a time counter and a packed state;
  \item a final theorem may expose only a code for the original input and a requested precision.
\end{itemize}
Forcing all lemmas to be stated from the beginning in the largest global frame is possible but makes every
formula unreadable and every binder manipulation fragile.

Let a local theorem be
\[
  d:\quad \Delta(u_0,\ldots,u_{k-1})\vdash A(u_0,\ldots,u_{k-1}).
\]
Let the final theorem have global variables \(x_0,\ldots,x_{n-1}\), and let
\(\sigma=(t_0(\vec x),\ldots,t_{k-1}(\vec x))\) be a substitution of global terms for the local variables.
The checked substitution theorem gives
\[
  d[\sigma]:\quad \Delta[\sigma]\vdash A[\sigma].
\]
If every formula of \(\Delta[\sigma]\) is available in a larger current context \(\Theta\), weakening gives
\(\Theta\vdash A[\sigma]\).

\begin{lemma}[Instantiation and context alignment]
\label{lem:instantiate-into}
For every local object derivation \(d:\Delta\vdash A\), term substitution \(\sigma\), and context
inclusion \(\Delta[\sigma]\subseteq\Theta\), there is an object derivation
\[
  \Theta\vdash A[\sigma].
\]
\end{lemma}
\begin{proof}
Apply \texttt{Prf.subst} to \(d\), obtaining
\(\Delta[\sigma]\vdash A[\sigma]\).  Apply \texttt{Prf.weaken} with the supplied list-inclusion proof.
No induction on formulas or derivations is repeated here; both operations are already checked generically.
\end{proof}

The supplied Lean file names this theorem \texttt{instantiateInto}.  Its code is essentially one line:
\begin{lstlisting}[language={}]
(local.subst sigma).weaken available
\end{lstlisting}
This compactness is important.  The trusted compiler kernel should be boring.

\subsection{Proof-carrying schedules}

Fix a permanent base context \(\Gamma\) over the global arity \(n\).  A reverse-topological proof schedule
with completed goals \(D=[A_{r-1},\ldots,A_0]\) carries the invariant
\[
  \forall A\in D,\quad \Gamma\vdash A.
\]
This is the structure \texttt{ProofBuild Gamma done} in the Lean outline.

Suppose a new local proof has already been aligned to
\[
  D+\!+\Gamma\vdash B.
\]
For every hypothesis of this derivation, either it is a prior goal in \(D\), in which case the schedule
contains a proof from \(\Gamma\), or it is a base formula in \(\Gamma\), in which case the hypothesis rule
proves it.  Simultaneous cut replaces all these hypotheses at once and produces \(\Gamma\vdash B\).
The schedule extends to \(B::D\).

\begin{lemma}[One schedule step]
If every formula in \(D\) is derivable from \(\Gamma\) and
\(D+\!+\Gamma\vdash B\), then every formula in \(B::D\) is derivable from \(\Gamma\).
\end{lemma}
\begin{proof}
For the new head \(B\), apply \texttt{Prf.trans} to the local derivation.  Map a hypothesis in \(D\) to the
stored schedule proof and a hypothesis in \(\Gamma\) to \texttt{Prf.hyp}.  Old members use the original
schedule invariant.
\end{proof}

An application node now has exactly the following data:
\begin{enumerate}
  \item a local arity \(k\);
  \item a local context \(\Delta\) and local goal \(A\) over \(k\) variables;
  \item a substitution \(\sigma:\mathrm{Fin}(k)\to\PTm(n)\);
  \item a local proof \(\Delta\vdash A\);
  \item an availability proof \(\Delta[\sigma]\subseteq D+\!+\Gamma\).
\end{enumerate}
The compiler applies Lemma~\ref{lem:instantiate-into} and then one schedule step.  All domain-specific
content is in the local proof and availability certificate; all structural content is in the generic
kernel.

\subsection{The static graph theorem}

\begin{definition}[Static proof graph]
A static proof graph over \(\Gamma\) is a finite sequence of nodes \(v=0,\ldots,m-1\).  Node \(v\) has
local data \((k_v,\Delta_v,A_v,\sigma_v,d_v)\).  Its availability condition says that every formula in
\(\Delta_v[\sigma_v]\) is either in \(\Gamma\) or equals \(A_u[\sigma_u]\) for some \(u<v\).
\end{definition}

\begin{theorem}[Staged graph compilation]
\label{thm:staged-graph}
For every static proof graph over \(\Gamma\), and every node \(v<m\),
\[
  \Gamma\vdash A_v[\sigma_v].
\]
In particular, the designated final node yields an object-iPRA derivation of the graph's public theorem.
\end{theorem}

\begin{proof}
Induct externally on the finite node sequence.  The empty sequence has the vacuous schedule.  At node
\(v\), the induction hypothesis is the proof-carrying schedule for all earlier nodes.  Apply
Lemma~\ref{lem:instantiate-into} to \(d_v\) and \(\sigma_v\), weakening to the concatenation of earlier
node goals and \(\Gamma\) using the availability certificate.  Apply one schedule step, whose simultaneous
cut discharges every earlier goal by the induction-hypothesis schedule and every base formula by a
hypothesis proof.  The extended schedule proves the claim for all nodes through \(v\).
\end{proof}

This induction occurs in the Lean metatheory while constructing a finite proof term.  It is not an
induction axiom added to iPRA, and the graph is not interpreted by an arithmetic universal machine.  Once
the finite graph is fixed, Lean elaboration produces a nested term made from substitution, weakening, cut,
and the local proof terms.

\subsection{Why topological order is the right certificate}

A graph description may use symbolic names and dependency edges.  The host compiler first checks that all
dependencies point backward in a supplied ordering.  This check is finite and decidable.  It then produces
an availability proof for each node.  There are two implementation styles.

\begin{enumerate}
  \item \emph{Typed construction.}  Build the graph incrementally with a dependent type whose constructor
        accepts only dependencies in the current list.  Ill-ordered graphs cannot be represented.
  \item \emph{Validated data.}  Parse a simpler JSON or Lean structure, run a topological sort, and use the
        result to generate a sequence of \texttt{NodeObligation} terms.  The generated Lean file contains
        the concrete list-inclusion proofs, so the kernel does not trust the parser.
\end{enumerate}
The second style is more convenient for large applications.  The final trust boundary remains Lean's
kernel checking the generated source.

\subsection{Recursive nodes}

Acyclic composition alone does not express a bounded iteration.  A loop node is compiled locally before it
is inserted into the graph.  Let \(P\) be a precondition, \(Q\) a postcondition, \(I(t,s,\vec x)\) a
quantifier-free invariant, and \(F(\vec x)\) a fuel term.  The local obligations are
\begin{align*}
  &P,\Gamma\vdash I(0,s_0,\vec x),\\
  &I(t,s_t,\vec x),(P,\Gamma)^{\uparrow}
        \vdash I(St,s_{St},\vec x),\\
  &I(F,s_F,\vec x),P,\Gamma\vdash Q.
\end{align*}
The shift on the permanent context is the ordinary binder shift required by universal introduction.  The
checked \texttt{PRA.qfInd} produces \(P,\Gamma\vdash\forall t\,I(t,s_t,\vec x)\); universal elimination at
\(F\) gives the final invariant; one cut with the finish proof yields \(P,\Gamma\vdash Q\), and implication
introduction gives \(\Gamma\vdash P\to Q\).

\begin{theorem}[Loop compilation]
\label{thm:loop-compile}
Every loop package with a quantifier-free invariant, base proof, step proof, and finish proof compiles to an
actual object derivation \(\Gamma\vdash P\to Q\).
\end{theorem}
\begin{proof}
The preceding paragraph is exactly the proof.  The supplied \texttt{LoopObligation.compile} follows this
sequence using \texttt{PRA.qfInd}, \texttt{Prf.allE}, and \texttt{Prf.trans}.
\end{proof}

A loop contract is then just another local theorem for the static graph.  Nested loops are handled by first
compiling the inner loop to a local implication theorem and using it in the outer step proof.  This
hierarchical strategy keeps each invariant small.  A compiler may flatten several loops into one packed
state, but flattening is an optimization, not a logical requirement.

\subsection{Substitution under binders and the reason to reuse the checked library}

Many earlier outlines become vague precisely at quantifier boundaries.  If a local theorem contains a
universal or existential formula, substituting global terms for its free variables must lift the
substitution under the binder.  The checked \texttt{Prf.substWeaken} proof already handles this.  In its
universal-introduction case it uses the lifted substitution \texttt{up sigma}; in its existential-
elimination case it aligns the shifted contexts and proves the relevant substitution identity.  Reusing
this theorem avoids a large class of de Bruijn-index errors.

The staged compiler therefore never manually rewrites a proof under a binder.  It asks
\texttt{Prf.subst} to specialize the entire derivation.  Domain libraries may state lemmas in their natural
local arity, and the compiler handles all binders uniformly.

\subsection{Native export is the final step}

After graph compilation yields
\[
  h:\Gamma\vdash_{\PRAAx}A,
\]
the checked bridge gives
\[
 \texttt{pra\_iff\_native.mp}(h):
 \NPrf(\Gamma^{\mathrm{native}},A^{\mathrm{native}}).
\]
Developing the graph directly in the very small native syntax would make application proofs needlessly
opaque.  Ordinary iPRA has convenient logical connectives and quantifiers; the bi-interpretation is already
proved.  The correct workflow is therefore: formalize the mathematics in ordinary \(\Prf\,\PRAAx\), audit
the endpoint, and export once.

\subsection{Trusted-code and proof-size accounting}

The generic logical kernel consists of:
\begin{itemize}
  \item one specialization-and-alignment theorem;
  \item the proof-schedule invariant;
  \item one simultaneous-cut step;
  \item one quantifier-free loop compiler;
  \item small wrappers for witnesses, stable outputs, boundary assumptions, and gap-based exactness.
\end{itemize}
All heavy programs and formulas are data supplied by application modules.  A graph with \(m\) nodes
produces a proof term with \(O(m+e)\) structural combinators, where \(e\) is the total number of dependency
uses, plus the sizes of the local proofs.  Repeatedly expanding a prior proof can create a tree larger than
this textual bound, but Lean's term sharing and named theorem constants avoid physical duplication.
The Python companion's \texttt{ProofTerm} trees are explanatory evidence only; the Lean kernel checks the
actual generated terms.
